\section{Conclusion and Limitations}





We benchmarked 11 deep learning optimizers in pretraining and found that their true gains over AdamW are much smaller than previously reported. Our results highlight three key lessons: (i) many claimed speedups stem from insufficient hyperparameter tuning, as fair sweeps eliminate most of the apparent advantages; (ii) comparisons based on early or inconsistent evaluation can be misleading, since optimizer rankings often change over the full training trajectory; and (iii) even the best-performing alternatives provide only modest speedups, which further diminish with model scale, dropping to ~1.1× at 1.2B parameters. This benchmarking study has the limitation that it does we haven't scaled to models larger than 1.2B parameters. However, we believe that evaluating optimizers on models of comparable size to prior studies is still valuable, as it reveals that insufficient tuning is a major cause of subpar speedups. Promising future work includes extending our benchmarking to larger models beyond 1.2B parameters to test whether the diminishing speedup trend persists at frontier scales. Another direction is to design optimizers whose efficiency remains stable under scaling laws, ensuring consistent benefits as model size and data budget grow.


\section*{Acknowledgement}

This work was supported by the Google TPU Research Cloud (TRC), the Stanford HAI–Google Cloud Credits Program, and  NSF RI 2045685, and is a part of the Marin Project. The authors thank Evan Walters, Omead Pooladzandi, Jiacheng You, and Zhiyuan Li for discussions and help during our research.